% ID: SEL-D0019-OA-E02
% ESTADO: validada. S020, cotejada con examen y criterios oficiales.
\textbf{a) Derivadas.} Por la regla del producto, para $x>0$,
$f'(x)=\ln x+x/x=\boxed{\ln x+1}$.
Por la regla del cociente, para todo $x\in\mathbb R$,
\[
g'(x)=\frac{3e^{3x}(x^4+1)-4x^3e^{3x}}{(x^4+1)^2}
=\boxed{\frac{e^{3x}(3x^4-4x^3+3)}{(x^4+1)^2}}.
\]

\textbf{b) Tangente.} Como $h(-2)=-3$ y $h'(-2)=2$, la recta
por $(-2,-3)$ de pendiente $2$ es $\boxed{y=2x+1}$.
La parábola $h(x)=(x+3)^2-4$ tiene vértice $(-3,-4)$ y corta al eje
horizontal en $-5$ y $-1$. La diferencia $h(x)-(2x+1)=(x+2)^2$
es cero justo en el punto de tangencia.
\begin{center}
\setlength{\unitlength}{6mm}
\begin{picture}(10,12)(-8,-8)
\qbezier(-7,0)(-3,0)(1,0)\put(1,0){\makebox(0,0)[l]{$x$}}
\qbezier(0,-7)(0,-1.5)(0,4)\put(0,4){\makebox(0,0)[b]{$y$}}
\qbezier(-5.5,2.25)(-3,-10.25)(-0.5,2.25)
\qbezier(-4,-7)(-2,-3)(0,1)
\put(-2,-3){\makebox(0,0){$\bullet$}}\put(-1.85,-3.35){\makebox(0,0)[l]{$(-2,-3)$}}
\put(-5.7,2.25){\makebox(0,0)[r]{$h$}}
\put(-4.1,-7){\makebox(0,0)[r]{$y=2x+1$}}
\end{picture}
\end{center}
